\documentclass[11pt]{article}
\usepackage[margin=0.85in]{geometry}
\usepackage{amsmath,booktabs,hyperref}
\title{Why the turning-crossing controller fails: controlled frame ablations}
\author{Pan Dynamics Collision Avoidance Research}
\date{October 2, 2026}
\begin{document}
\maketitle
\section{Finding}
The 0.85 s frame contains two distinct approximation failures. The fixed
separation directions can exclude collision-free turning configurations;
enlarging the input box then finds affine zero-slack plans whose nonlinear
motion collides. The prior explanation in terms of excessive linearization
error alone was incomplete. A controlled normal perturbation isolates the
geometric restriction, while unchanged dynamics and primary/secondary
comparisons isolate the misleading safety objective and its interaction with
CLF improvement. No production controller modification was made.

All experiments start from the recorded state of the final 15 m/s
turning-crossing case at 0.85 s. The exact saved fresh-flow anchor and affine
problem are reused. The original sample time is 0.05 s, with 16 primary and
60 completion intervals, a 0.10 m buffer, exact observation, known target
motion, and no road constraints. The original numerical input correction
box is $\pm0.075$ rad and $\pm0.125$ signed longitudinal ratio. The source
under diagnosis is commit
\begin{center}\small\texttt{2697350ca383bbbc326f2ca516129c740cc29fd0}.\end{center}

\section{The geometric restriction is directly observable}
Actual rectangle separation can be expressed by existence of a separating
normal $n$. The implementation selects $n$ from each anchor configuration
and freezes it within that convex problem. Normals may differ between
prediction nodes and are recomputed on a new anchor, but do not adapt to the
optimized pose within the two-stage solve. Vertex yaw dependence is
linearized with that normal fixed.

Using the returned inputs from a normal-perturbed problem, propagate the
nonlinear vehicle and inspect the same actual configuration at 1.65 s:
\begin{center}
\begin{tabular}{lr}
\toprule
Quantity & Value\\
\midrule
Actual rectangle distance & 0.105235 m\\
Projected gap using old normal & $-0.160752$ m\\
Old normal angle & $112.84^\circ$\\
Normal consistent with actual separation & $116.02^\circ$\\
\bottomrule
\end{tabular}
\end{center}
The vehicle poses are identical in the two geometry calculations. The old
supporting line cuts through a rear corner of the ego even though that corner
is far from the target. The normal consistent with the turned pose separates
the bodies. Thus failure of this particular fixed-normal certificate is not
proof of actual overlap. This conservatism exists even with exact vehicle
states and exact rectangle support calculations.

The earlier one-row LP bound remains valid for the original local problem:
a hard corner row at stage 17 needs 0.390161 m of improvement but has at most
0.342086 m available under the affine dynamics and variable bounds. Its
best projected gap is 0.051924 m, short of the required 0.10 m. All other
inequality rows and the endpoint cone were omitted in that bound. The new
experiment identifies that this impossibility is tied to the chosen direction
as well as the allowed corrections; it is not a universal physical limitation.

\section{Controlled optimization experiments}
Every reconstructed start/midpoint safety row is checked against the saved
problem before perturbation. Normal experiments rotate all such directions
with timestamps in $[1.4,1.9]$ s, preserving the anchor, dynamics, bounds,
endpoint cone and primary/tail structure. The following are raw primary
results; their complete input sequences are propagated nonlinearly with
5-ms RK4 substeps through 4.65 s.
\begin{center}\small
\begin{tabular}{lrrr}
\toprule
Change from original local problem & Exit flag & PCBF sum & Min. distance [m]\\
\midrule
None & $-2$ & -- & --\\
Normal rotated $-20,-10,-5^\circ$ & $-2$ & -- & --\\
Normal rotated $+5^\circ$ & 1 & 0.158620 & 0.094136\\
Normal rotated $+10^\circ$ & 1 & 0.586658 & 0.094136\\
Normal rotated $+20^\circ$ & 1 & 2.076820 & 0 (overlap)\\
Buffer reduced to zero & 1 & $6.06\,10^{-9}$ & 0 (overlap)\\
Input box enlarged $1.25\times$ & 1 & $1.07\,10^{-8}$ & 0 (overlap)\\
Input box enlarged $1.5\times$ & 1 & $3.58\,10^{-9}$ & 0 (overlap)\\
Input box enlarged $2\times$ & 1 & $3.04\,10^{-7}$ & 0 (overlap)\\
State trust removed, physical domain retained & $-2$ & -- & --\\
\bottomrule
\end{tabular}
\end{center}
Positive primary slack and positive actual clearance are compatible: fixed
supporting directions are only sufficient tests, and the surrogate includes a
0.10 m buffer. The 9.4 cm minimum meets the requested noncollision criterion
at evaluated samples but falls slightly below that buffer. The table does not
claim zero-slack MPC feasibility, continuous-time safety, or terminal
membership for the nonlinear trajectory. Conversely, tiny affine slack in
the enlarged-box variants does not imply actual clearance.

These ablations establish that merely removing state trust, reducing the
buffer, or increasing the allowable input increment is not a reliable remedy
in this frame. The direction perturbation is diagnostic, not a proposed
fixed angular offset, alternative controller mode, or online normal bank.

\section{The CLF stage and nonlinear verification}
The unchanged secondary solver was applied to the $+5^\circ$, $+10^\circ$,
and original-direction $2\times$ cases. A temporary helper outside the
repository extracts the original secondary function and its local helpers;
only its entry-point name changes. It constructs the same CLF and enforces
the same primary objective cap.
\begin{center}\small
\begin{tabular}{lrrr}
\toprule
Variant & Returned safety sum & First $(\delta,b)$ & ODE45 minimum gap\\
\midrule
$+5^\circ$, original box & 0.158621 & $(0.203988,-0.474999)$ & 0.094135 m\\
$+10^\circ$, original box & 0.586659 & $(0.203987,-0.475000)$ & 0.094136 m\\
Original direction, $2\times$ box & $1.28\,10^{-6}$ & $(-0.021011,-0.100000)$ & 0 (overlap)\\
\bottomrule
\end{tabular}
\end{center}
Verification used ODE45 with relative tolerance $10^{-11}$, absolute tolerance
$10^{-12}$, and 11 output samples per hold. For the $+5^\circ$ result the
minimum occurs at 1.645 s. Speed remains within 10.147--15.008 m/s, maximum
lateral speed is 0.343 m/s and maximum yaw rate is 0.745 rad/s, within the
configured physical state domain. The CLF stage does not destroy the
noncollision behavior of this variant.

For the original enlarged problem, the primary first input is approximately
$(0.251855,-0.493934)$; the CLF first input becomes
$(-0.021011,-0.100000)$. The latter turns less favorably and brakes less while
still preserving almost-zero modeled safety slack. Both primary and secondary
nonlinear plans overlap. The earlier measured affine/nonlinear discrepancy
and extrapolated tire-force optimism therefore remain relevant. This is a
failure of the safety approximation that constrains CLF optimization, not
evidence that CLF priority was incorrectly ordered or skipped.

Twelve additional finite braking probes used zero steering and ratios
$-0.5,-0.65,-0.8$ for 0.4--1.0 s, followed by the existing nominal feedback
only for diagnostic propagation. All overlapped. No claim that simple
straight braking solves this encounter is supported. These probes are not
an online backup or an implemented policy.

\section{Causal interpretation and design implications}
The flow-guided seed overlaps the target at the start of the hard completion
part. It also changes its shaped longitudinal input from $-0.35$ to $+0.35$
before that conflict, so it is not itself a safe witness; an unsafe seed is
allowed by design. Its local normal selects a restrictive geometric
correction. With the original correction box the affine problem is
infeasible. Expanding the box yields an optimistic zero-slack affine plan,
then CLF optimization can favor controls whose actual nonlinear avoidance is
worse. At the next frame the imminent conflict moves from the hard completion
part into the slack-enabled primary part, after which positive safety slack
persists until the observed closed-loop collision.

The local slack values from different normals or anchors are consequently
not faithful rankings of nonlinear collision avoidance: here a positive-slack
variant avoids sampled overlap while a near-zero-slack variant collides.
The measurements support a redesign that coordinates separation geometry,
trajectory updates and model validity. They do not support treating a fixed
normal offset or an arbitrarily enlarged trust region as a complete fix.
Repeated relinearization must actually improve the coupled trajectory and
geometry; iteration count alone is not a convergence guarantee.

The noncolliding counterfactual shows an admissible-input vehicle motion over
this finite interval, not a full closed-loop scenario pass, a terminal-set
certificate, or an infinite-horizon guarantee. No general conclusion that
all nearby trajectories are feasible or that all failures have the same
cause is made.

\section{Artifacts and checks}
The dated companion directory contains three MATLAB diagnostic drivers,
a Python figure renderer, ablation tables, sampled counterfactual traces,
reproduction instructions and a manifest of original evidence. Generated
figures and temporary controller helpers remain outside the repository.
The first two panels of the figure use exactly the same vehicle poses and
show the old and pose-consistent normals. The third panel illustrates the
prior one-row correction bound. Independent Python rectangle geometry
reproduces the 0.105235 m distance. Source-row reconstruction assertions,
ODE45 comparisons, physical-domain summaries, artifact hashes, report
compilation and staged whitespace checks were performed. No full scenario
campaign or unit suite was rerun, since production code is unchanged.
\end{document}
